\section{Tło projektu}

Rynek motoryzacyjny jest jednym z kluczowych sektorów, wywierających wpływ na polską gospodarkę, a jego znaczenie rośnie co roku.
W samym 2025 liczba zarejestrowanych samochodów osobowych w Polsce wyniosła 4\,249\,151 sztuk, czyli o 3\% więcej niż w roku 2024 \autocite{otomoto_2025}.
Istotną część polskiej branży motoryzacyjnej pełni rynek wtórny, na którym sprzedano łącznie 3\,345\,928 używanych samochodów w 2025 roku \figref{fig:automotivemarket}.
Średni wiek pojazdu osobowego sprzedanego na takim rynku wynosił 11 lat, a średni przebieg ponad 168 tysięcy kilometrów.
Rynek ten jest bardzo dynamiczny, ponieważ średni czas sprzedaży samochodu osobowego zajmował około 40 dni \autocite{aaaauto_2025}.
%
\begin{figure}[h]
    \centering
    \includegraphics[width=0.8\textwidth]{automotive-secondary-market_2025}
    \figcaption{Wysokość sprzedaży oraz średniej ceny samochodów na polskim rynku wtórnym w 2025 roku}{\cite{aaaauto_2025}}
    \label{fig:automotivemarket}
\end{figure}

W celu znalezienia auta nadającego się do zakupu, Polacy często wybierają serwisy internetowe z ogłoszeniami o sprzedaży samochodów.
Przykładowo serwis OtoMoto odnotował, w grudniu 2025 roku, ponad 11 milionów unikalnych odwiedzin witryny \autocite{otomoto_2025}.
Pomimo popularności tych serwisów, większość kupujących wyraża chęć osobistej inspekcji auta przed zakupem.
Badanie przeprowadzone przez \emph{carVertical} wskazuje na to, że ponad 80\% polskich respondentów przyznało, że nie ufa sprzedawcom.
Wynika to często z negatywnych doświadczeń, w których pojazdy miały cofane liczniki przebytych kilometrów lub zatajone szkody \autocite{motofaktor_trust}.

\section{Istniejące serwisy}

\subsection{OtoMoto}

Serwis \emph{OtoMoto} jest liderem w sektorze kupna i sprzedaży pojazdów, zarówno nowych jak i używanych.
Ogłoszenia tam mogą wystawiać osoby prywatne, komisy oraz autoryzowani dealerzy marek.
Serwis cechuje się zaawansowaną wyszukiwarką z licznymi filtrami, które pozwalają dokładnie przeglądać dostępne ogłoszenia.
Dodatkowym atutem jest fakt, że ogłoszenia mogą zostać oznaczone jako zweryfikowane pod względem zgodności danych ze stanem faktycznym.
Większość ogłoszeń jednak jej nie podlega, przez co kupujący często spotykają się z ogłoszeniami, które wykorzystują nieuczciwe praktyki.
Ponadto serwis w wynikach wyszukiwania wyróżnia oferty, które nie są w pełni zgodne z filtrami, ale zostały odpowiednio opłacone przez podmioty je wystawiające.
Sam wymóg uiszczenia opłaty za listowanie ogłoszenia na serwisie, jest wątpliwe w skuteczności zapobiegania nieuczciwym praktykom.
Użytkownicy nie mają narzędzi do przeciwstawiania się takim ofertom innych niż zgłoszenie ogłoszenia, które posiada wątpliwą skuteczność.
Sprzedający są anonimowi, a użytkownicy nie mają dostępu do historii wystawianych przez nich ofert.

\subsection{OLX}

Jednym z najpopularniejszych serwisów do wystawiania ogłoszeń o sprzedaży jest portal \emph{OLX}.
Nie jest to w pełni skoncentrowany na motoryzacji serwis, natomiast są na nim kategorie obejmujące sprzedaż pojazdów mechanicznych.
Serwis posiada rozbudowany zakres filtrów dla ogłoszeń, w tym zakładki pozwalające podzielić wyniki wyszukiwania na prywatnych sprzedawców oraz firmy.
Możliwości weryfikacji ogłoszeń są ograniczone do zera, a zespół moderujący treści na stronie jest nieskuteczny w usuwaniu nieuczciwych ofert.
Serwis również stosuje wyróżnianie wyników, zależne od wniesionych przez sprzedawców opłaty, jednak stawki są niezależne od wartości sprzedawanego pojazdu i niższe niż w przypadku OtoMoto.
Sprzedawcy posiadają profile, w których można sprawdzić aktualne oferty, ale ich archiwum nie jest udostępniane.
Serwis natomiast posiada możliwość wystawiania ocen, jednak tylko po zakończeniu atrakcji, co jest znacząco niewystarczające w kontekście motoryzacyjnym.
Problemem serwisu jest również fakt, że udostępnia on ogłoszenia z portalu OtoMoto, przez co przeglądanie wyników wyszukiwania jest wolniejsze ze względu na przenoszenie do zewnętrznego serwisu.
Ponadto, takie działanie skutkuje tym, że kupujący widzą to samo ogłoszenie ponownie, co utrudnia szukanie odpowiedniej oferty.

\subsection{Allegro}

Popularną internetową platformą handlową jest serwis \emph{Allegro}, który posiada specjalną kategorię dla ogłoszeń motoryzacyjnych.
W przeciwieństwie do innych produktów, platforma nie pośredniczy w transakcjach pomiędzy sprzedawcą i kupującym, a udostępnia jedynie możliwość publikowania ofert.
Z tego też powodu, działający na serwisie system ochrony kupujących nie obejmuje zakupu pojazdów i wsparcia w przypadku oszustwa.
Podobnie jak w przypadku innych serwisów, wystawianie ogłoszeń jest odpłatne z dodatkowymi ofertami promocji.
Serwis udostępnia wiele filtrów dla manipulacji wynikami wyszukiwania zgodnych z innymi kategoriami, natomiast w przypadku motoryzacji są bardzo niewygodne w użyciu.
Zespół moderujący nie weryfikuje ogłoszeń pod względem ich zgodności ze stanem faktycznym, natomiast występują tam rozbudowane narzędzia do tworzenia zgłoszeń.
Analogicznie jak na serwisie OLX, kupujący mogą wystawiać opinie dopiero po zakupie, co nie sprawdza się dla samochodów osobowych.

\section{Cel i zakres pracy}

Celem niniejszej pracy jest wytworzenie rozwiązania, które skupiałoby się na poprawieniu obecnej sytuacji kupujących, na rynku wtórnym sprzedaży pojazdów mechanicznych.
Wymienione serwisy nie posiadają wystarczających narzędzi, aby kupujący mogli ochronić się przed nieuczciwymi praktykami sprzedających.
Równocześnie istniejące systemy moderowania treści nie są w stanie zapewnić im wymaganego bezpieczeństwa.

Pierwotnym założeniem projektu było stworzenie nakładki, rozszerzającej funkcjonalności obecnych serwisów o system wystawiania opinii z poziomu odrębnej platformy.
Klienci docelowi mogliby wyszukiwać oferty z poziomu takiej nakładki, połączonej interfejsem programowania aplikacji z portalami ogłoszeń o sprzedaży samochodów, a następnie sprawdzać opinie innych użytkowników na ich temat.
Jeżeli społeczność serwisu nie miałaby zastrzeżeń co do wystawionej oferty, wtedy można by było wybrać dane ogłoszenie i zostać przekierowanym na właściwy serwis, na którym zostało ono wystawione.
Jednak serwisy te nie wystawiają swoich interfejsów publicznie do sieci, a dostęp do nich albo nie jest darmowy albo w ogóle nie udostępniany dla prywatnych podmiotów, więc projekt w tej postaci nie mógł zostać zrealizowany.

Niniejsza praca polega na stworzeniu serwisu z ogłoszeniami o sprzedaży samochodów, który będzie posiadał narzędzia ułatwiające przeciwdziałanie oszustom.
Realizacja tego rozwiązania opiera się na spełnieniu poniższych wymagań funkcjonalnych:
%
\begin{itemize}
    \item możliwość zamieszczania komentarzy oraz ocen pod ofertami,
    \item system reputacji bazujący na wszystkich dotychczasowych ofertach sprzedawcy oraz pozycjonowanie ogłoszeń na jego podstawie,
    \item archiwum ogłoszeń danego sprzedającego, które jest publicznie dostępne,
    \item pozycjonowanie opinii na podstawie czynników takich jak, czy użytkownik był na oględzinach pojazdu,
    \item skuteczna wyszukiwarka ofert, które interesują użytkowników,
    \item system weryfikowania, moderowania oraz zgłaszania ofert oraz opinii.
\end{itemize}
